\subsection{Introduction}
x86lite is a subset of the full x86 ISA, with only about 20 instructions and only 64-bit signed integers, nothing else.
It is further interoperable with x86.

As a quick reminder, the following registers are relevant for x86lite:
\begin{multicols}{4}
    \begin{itemize}
        \item \texttt{rax}
        \item \texttt{rbx}
        \item \texttt{rcx}
        \item \texttt{rdx}
        \item \texttt{rsi}
        \item \texttt{rdi}
        \item \texttt{rbp} (Base Pointer)
        \item \texttt{rsp} (Stack Pointer)
        \item \texttt{r08}
        \item \texttt{r09}
        \item \texttt{r10}
        \item \texttt{r11}
        \item \texttt{r12}
        \item \texttt{r13}
        \item \texttt{r14}
        \item \texttt{r16}
    \end{itemize}
\end{multicols}

Further, the \texttt{rip} points to the next instruction. Then, the CPU has processor state registers, such as \texttt{OF} (set, if the last instruction caused an overflow),
\texttt{ZF} (set if last result was zero), etc.

On x86, the heap grows upwards, the stack grows downwards, where the code and data section that precedes the heap in the address space contains program code, constants and globals.
